\documentclass[12pt]{article}
\usepackage[utf8]{inputenc}
\usepackage[T1]{fontenc}
\usepackage[margin=1in]{geometry}
\usepackage{setspace}
\usepackage{xcolor}
\usepackage[colorlinks=true, urlcolor=blue!60!black, linkcolor=blue!60!black]{hyperref}
\setlength{\parindent}{2em}
\setlength{\parskip}{1.25em}
\renewcommand{\baselinestretch}{1.0}
\begin{document}
\begin{center}
{\LARGE \textit{National Pork Producers Council v. Ross}, 598 U.S. 356 (2023)\par}
\vspace{0.8em}
{\large Case Brief --- Four-Page Edition\par}
\end{center}
\vspace{1em}

\section{Memory Jogger}

Affirming dismissal of a Dormant Commerce Clause challenge to California's Proposition 12, the Court unanimously rejected an ``almost per se'' rule against state laws with extraterritorial effects and fractured 5--4 in upholding California's nondiscriminatory ban on in-state sales of pork from cruelly confined breeding pigs, leaving uniform national rules to Congress (\href{https://www.law.cornell.edu/supremecourt/text/21-468}{598 U.S.\ 356}).

\section{Detailed Case Facts}

In 2018, about 63\% of California voters approved Proposition 12, which barred the in-state sale of whole pork meat from a breeding pig, or its immediate offspring, ``confined in a cruel manner,'' meaning unable to lie down, stand up, extend its limbs, or turn around freely, with minimum space standards.

The National Pork Producers Council and the American Farm Bureau Federation sued California officials; the Humane Society intervened. They alleged that California consumes about 13\% of the nation's pork but imports nearly all of it, and that compliance would cost \$290 to \$348 million and raise production costs 9.2\%, practically forcing producers nationwide to restructure. They conceded that the law treats in-state and out-of-state producers alike, that 28\% of the industry already used group housing, and that Congress had declined to preempt such laws.

\section{Procedural History}

The Southern District of California dismissed under Rule 12(b)(6). A unanimous Ninth Circuit panel (Ikuta, J.) affirmed, holding that upstream effects and compliance costs neither violated the extraterritoriality doctrine nor established a substantial burden under \textit{Pike v.\ Bruce Church}. The Supreme Court granted certiorari and affirmed on May 11, 2023.

\section{Judicial Votes}

Affirmed, \textbf{5--4}.
\begin{itemize}
\item \textbf{Judgment and majority (Parts I, II, III, IV--A, V):} Justice Gorsuch, joined by Justices Thomas, Sotomayor, Kagan, and Barrett. All nine Justices rejected the ``almost per se'' rule.
\item \textbf{Part IV--B and IV--D (incommensurability):} a three-Justice plurality (Gorsuch, Thomas, Barrett); six Justices rejected it.
\item \textbf{Part IV--C (no substantial burden pleaded):} a four-Justice plurality (Gorsuch, Thomas, Sotomayor, Kagan); five Justices disagreed.
\item \textbf{Separate opinions:} Sotomayor (joined by Kagan) and Barrett concurring in part; Chief Justice Roberts (joined by Alito, Kavanaugh, and Jackson) and Kavanaugh concurring in part and dissenting in part.
\end{itemize}

\section{Holding}

The antidiscrimination principle lies at the ``very core'' of the Dormant Commerce Clause. There is no ``almost per se'' rule invalidating nondiscriminatory laws merely because they affect out-of-state commerce; \textit{Baldwin}, \textit{Brown-Forman}, and \textit{Healy} involved protectionist price controls. \textit{Pike} survives but chiefly smokes out hidden protectionism or protects instrumentalities of interstate transportation, and absent congressional preemption courts must use ``extreme caution'' before invalidating nondiscriminatory regulation of in-state sales. Because neither Part IV rationale won a majority, the controlling \textit{Pike} ground under \textit{Marks} is contested.

\section{Analysis of Opinions}

\textbf{Gorsuch (majority and pluralities).} Congress could enact a uniform pork standard but has not. Petitioners' rule would imperil countless ordinary state laws. \textit{Pike} itself smoked out protectionism, and non-protectionist invalidations rarely go beyond transportation: ``Pigs are not trucks or trains.'' For three Justices, weighing producers' costs against voters' moral interests is like deciding ``whether a particular line is longer than a particular rock is heavy,'' a choice for voters and Congress. For four, petitioners pleaded no substantial burden under \textit{Exxon}, because the Clause protects the interstate market, not particular firms' ``methods of operation.''

\textbf{Sotomayor (joined by Kagan).} \textit{Pike} balancing remains workable and reaches beyond discrimination and transportation; she affirmed only because no substantial burden was pleaded.

\textbf{Barrett.} The complaint did allege substantial burdens, but California's moral interest cannot be weighed on a judicial scale against dollars, so she affirmed.

\textbf{Roberts (joined by Alito, Kavanaugh, Jackson).} Six Justices preserved \textit{Pike} and five found a burden plausibly alleged; under \textit{Bibb}, \textit{Kassel}, and \textit{Edgar}, market-wide harms to farmers nationwide count. He would have remanded for balancing.

\textbf{Kavanaugh.} California's ``California knows best'' regulation threatens balkanization; he flagged the Import-Export, Privileges and Immunities, and Full Faith and Credit Clauses for future challenges.

\section{Examples: Future Applications}

\begin{enumerate}
\item \textbf{Same-Side (Upheld): Compostable Packaging Mandate.} A nondiscriminatory in-state sales rule requiring compostable packaging only changes firms' methods of operation, absent congressional preemption.
\item \textbf{Same-Side (Upheld): Fur or Foie Gras Sales Ban.} A State may ban in-state sales of goods on nondiscriminatory moral or animal-welfare grounds.
\item \textbf{Opposite-Side (Invalid): Lowest-Price Matching.} A law tying in-state drug prices to the lowest neighboring-State price falls under the preserved \textit{Baldwin}--\textit{Healy} price-control line.
\item \textbf{Opposite-Side (Invalid): Unique Railcar Couplers.} Conflicting coupler mandates that force trains to reassemble at the border burden ``instrumentalities of interstate transportation.''
\item \textbf{Fence-Sitter: Sales Bans Tied to Out-of-State Labor or Reproductive Policies.} Facially neutral bans tied to unrelated corporate policies escape the per se rule but invite \textit{Pike} claims, given five Justices' view of burden, and the clauses Kavanaugh flagged.
\end{enumerate}

\section{Critique}

\textbf{Originalist and textualist views.} Originalists praised Gorsuch's refusal to turn the atextual Dormant Commerce Clause into a ``roving license'' for judges, leaving national rules to Congress. Some free-market commentators worry about balkanization by dominant States and welcomed Kavanaugh's invitation to revisit \textit{Woodruff v.\ Parham}.

\textbf{Progressive and doctrinal views.} Progressive and animal-welfare scholars hailed a win for state autonomy and democratic moral choice. Doctrinal commentators note that the voting paradox leaves \textit{Pike} uncertain and, as the \textit{Harvard Law Review} observed, that deference favors large States able to export their rules nationwide.

\section{Key Quotations}

\begin{quote}
``Companies that choose to sell products in various States must normally comply with the laws of those various States. Assuredly, under this Court's dormant Commerce Clause decisions, no State may use its laws to discriminate purposefully against out-of-state economic interests. But the pork producers do not suggest that California's law offends this principle. Instead, they invite us to fashion two new and more aggressive constitutional restrictions on the ability of States to regulate goods sold within their borders. We decline that invitation. While the Constitution addresses many weighty issues, the type of pork chops California merchants may sell is not on that list.''
\end{quote}
\noindent --- Justice Gorsuch, Opinion of the Court, 598 U.S.\ at 364.

\begin{quote}
``How is a court supposed to compare or weigh economic costs (to some) against noneconomic benefits (to others)? No neutral legal rule guides the way. The competing goods before us are insusceptible to resolution by reference to any juridical principle. Really, the task is like being asked to decide `whether a particular line is longer than a particular rock is heavy.' ''
\end{quote}
\noindent --- Justice Gorsuch, Plurality Opinion, Part IV--B, 598 U.S.\ at 381.

\begin{quote}\sloppy
``The State has aggressively propounded a `California knows best' economic philosophy---where California in effect seeks to regulate pig farming and pork production in all of the United States. California's approach undermines federalism and the authority of individual States by forcing individuals and businesses in one State to conduct their farming, manufacturing, and production practices in a manner required by the laws of a different State.''
\end{quote}
\noindent --- Justice Kavanaugh, Concurring in Part and Dissenting in Part, 598 U.S.\ at 406--07.

\end{document}
